\section{Estrutura do trabalho}
Este estudo está estruturado em mais cinco capítulos distintos. No capítulo \ref{sec:estado_atual}, apresenta-se em detalhes a metodologia abrangente que norteou o desenvolvimento do Vortex, abordando desde sua fundação tecnológica até suas políticas, padrões de projeto e boas práticas de desenvolvimento seguidas. Nessa seção, são delineados os principais conceitos, funcionalidades e classes que compõem integralmente o referido \textit{framework}.

No capítulo \ref{sec:modelagem} estão descritas as classes e suas responsabilidades dentro do \textit{framework}, assim como suas relações com as demais, além das representações dos diagramas das mesmas.

O subsequente capítulo \ref{Utilização} fornece uma descrição pormenorizada do procedimento adequado para a utilização de cada funcionalidade do Vortex. Tal exposição abrange desde comandos do Cosmo até as classes que integram a arquitetura global do \textit{framework}.

O capítulo \ref{sec:estudo_de_caso} é dedicada ao estudo de caso, onde são delineados os testes conduzidos para verificar a superioridade de desempenho do Vortex em comparação com outros \textit{frameworks} similares. Essa comprovação é estabelecida mediante a realização de experimentos controlados em cenários equivalentes.

Finalmente, a conclusão, abordada no capítulo \ref{chp:conclusao}, destaca os resultados obtidos ao longo do estudo, proporcionando \textit{insights} conclusivos. Além disso, são delineadas as perspectivas para trabalhos futuros, consolidando assim a contribuição e relevância do presente trabalho.